% Generated draft: requires mathematical and source review.
\documentclass[11pt]{article}
\usepackage[margin=25mm]{geometry}
\usepackage{fontspec}
\setmainfont{Latin Modern Roman}
\usepackage{amsmath,amssymb,mathtools,unicode-math}
\setmathfont{Latin Modern Math}
\usepackage{enumitem,needspace,xurl,hyperref}
\setlength{\parindent}{0pt}
\setlength{\parskip}{5pt}
\setlength{\emergencystretch}{3em}
\newcommand{\N}{\mathbb{N}}
\newcommand{\Z}{\mathbb{Z}}
\newcommand{\Q}{\mathbb{Q}}
\newcommand{\R}{\mathbb{R}}
\newcommand{\C}{\mathbb{C}}
\newcommand{\F}{\mathbb{F}}
\newcommand{\A}{\mathbb{A}}
\newcommand{\PP}{\mathbb{P}}
\newcommand{\E}{\mathbb{E}}
\newcommand{\eps}{\varepsilon}
\newcommand{\dd}{\,\mathrm{d}}
\begin{document}
\section*{UnsolvedMath ID: 20002711}
\textbf{Source problem code:} catalogue entry
\subsection{Definitions and mathematical statement}
Let $W$ be a finite Coxeter group and $T$ its set of reflections. Write $\ell_T$ for reflection length, and define absolute order by
\[
 u\le_Tv\quad\Longleftrightarrow\quad
 \ell_T(v)=\ell_T(u)+\ell_T(u^{-1}v).
\]
Its rank is $\ell_T$, with minimum the identity $1$. V. Reiner's source question asks for its topology, EL-labelings/shellings in types $A$ and $B$, and whether passing to the ideal generated by parabolic Coxeter elements repairs the reported type-$D$ obstruction. A parabolic Coxeter element is a Coxeter element in a conjugate of a standard parabolic subgroup, with the intended Coxeter-system convention fixed. The remark asks whether Brady--Watt methods for noncrossing intervals generalize.

For a finite poset $P$, its order complex $\Delta(P)$ has finite chains as simplices. Literal $\Delta(\operatorname{Abs}(W))$ is a cone at $1$, hence contractible. The substantive convention in the aid is
\[
 K(W)=\Delta(\operatorname{Abs}(W)\setminus\{1\}).
\]
Deleting a minimum, deleting additional endpoints, and adjoining a maximum give different complexes. An EL-labeling labels cover edges in a total order so each interval has a unique increasing maximal chain, lexicographically first. The source does not fully fix its global labeling convention for this generally nonbounded poset.

The aid reports known homotopy Cohen--Macaulayness for the deleted type-$A$ complex and for the type-$B$ parabolic-Coxeter ideal, EL-shellability of noncrossing intervals, and a non-Cohen--Macaulay interval obstruction in $D_4$. Full hypotheses and references are absent; interval results do not automatically provide a global EL-labeling.

For the dihedral system $I_2(m)$, $m\ge3$, the aid supplies an exact family. There are $m$ reflections of rank one and $m-1$ nonidentity rotations of rank two, every reflection below every such rotation. Therefore
\[
 K(I_2(m))=K_{m,m-1}\simeq
 \bigvee^{(m-1)(m-2)}S^1,
\]
where $K_{a,b}$ denotes the complete bipartite graph. For the ideal generated by parabolic Coxeter elements, only the two full-rank Coxeter rotations remain above the reflections, yielding
\[
 \Delta(\mathcal J(I_2(m))\setminus\{1\})=K_{m,2}
 \simeq\bigvee^{m-1}S^1.
\]
The count uses standard parabolic Coxeter elements, not every arbitrary generating pair of reflections with a newly chosen Coxeter structure.

The aid also reports the product homotopy reduction
\[
 K(W_1\times W_2)\simeq K(W_1)*K(W_2),
\]
with the usual join/empty-complex conventions. In particular, for $q\ge1$ and $m_j\ge3$,
\[
 K\!\left(\prod_{j=1}^qI_2(m_j)\right)
 \simeq\bigvee^{\prod_j(m_j-1)(m_j-2)}S^{2q-1}.
\]
These exact rank-two and product statements do not settle the source's general type-$A,B$ EL-labeling or uniform parabolic-ideal program.
\subsection{Short English statement}
Keeping the identity makes the order complex contractible, so the aid studies the complex with that minimum deleted. It computes dihedral examples and a product reduction, while retaining the distinction between known interval shellings and the source's global topology and labeling questions.
\subsection{Sources}
\begin{itemize}
\item Local UnsolvedMath registry, permanent ID 20002711, source code catalogue entry; created and updated 14 August 2026. Canonical \texttt{aim-probability-notes.json}, zero-based \texttt{notes[152]}, tag \texttt{problem}, background item 3, attributed in the extract to V. Reiner.
\item AIM ``Braid groups, clusters and free probability'': \url{https://aimath.org/WWN/braidgroups/braidgroups.pdf}. The exact metadata label accompanies an extraction beginning ``Problem 3.'' followed by ``1.''; supplied verification says a section heading is merged and symbols damaged. No unverified renumbering is imposed here.
\item The original remark cites Brady and Watt [14] for noncrossing interval EL-labeling methods, without a full bibliography. The literature field is empty. Research attempt 1: \texttt{partial\_result}, problem status at run \texttt{partially\_solved}, rigor label \texttt{proved}, novelty confidence low; complete citations for the later topology assessments are absent.
\item Ulam AI / UnsolvedMath Contributors, CC BY 4.0. Machine-generated research aid, not peer reviewed; independently verify before citation.
\end{itemize}
\end{document}
